\documentclass[GREVA,english,openany]{tfe-eseiaat}

\usepackage{tikz}
\usetikzlibrary{shapes.geometric,shapes.symbols,positioning,calc}

\title{Study and Implementation of a Three-Dimensional Numerical Simulator Based on the Finite Element Method in Parallel Computing Environments}

\thisdocument{Project Charter}
\author{Marina-Liao Colom Escolá}
\email{marina-liao.colom@estudiantat.upc.edu}
\tutor{Rafel Amer Ramon}

\begin{document}

\chapter{Object}

The object of this project is the study and implementation of a numerical
simulator that solves boundary value problems posed on three-dimensional
domains and discretised with the finite element method, and that is able to
run in parallel on distributed-memory architectures.

Concretely, the simulator addresses the linear advection--diffusion--reaction
equation, both in its steady (elliptic) form
\begin{equation}
    -\nabla\cdot(a\,\nabla u) + b\cdot\nabla u + c\,u = f
    \qquad \text{in } \Omega \subset \mathbb{R}^{3},
    \label{eq:steady}
\end{equation}
and in its time-dependent (parabolic) form
\begin{equation}
    \frac{\partial u}{\partial t}
    -\nabla\cdot(a\,\nabla u) + b\cdot\nabla u + c\,u = f
    \qquad \text{in } \Omega \times (0, T],
    \label{eq:transient}
\end{equation}
where the diffusion coefficient $a$, the advection field $b$, the
reaction coefficient $c$ and the source term $f$ may all depend on the point
$(x, y, z)$, and equation~\eqref{eq:steady} is closed with Dirichlet, Neumann
or Robin boundary conditions on the different parts of the boundary
$\partial\Omega$.

The implementation is built on top of PETSc, the \emph{Portable, Extensible
Toolkit for Scientific Computation} \cite{petsc-user-ref}, which provides the
distributed data structures and parallel solvers on which the simulator
relies. Unstructured tetrahedral meshes of the three-dimensional domains are
generated externally with Gmsh \cite{gmsh-manual}, and the definition of each
problem ---its coefficients, source terms and boundary conditions--- is
supplied at run time through an embedded Lua scripting layer, so that a new
problem can be solved without recompiling the code.

\begin{figure}[!ht]
\begin{center}
\includegraphics[width=0.8\textwidth]{meshes/torus-int}
\end{center}
\caption{Unstructured tetrahedral mesh of a solid torus.}
\label{fig:mesh3d}
\end{figure}

This work also builds upon, and extends to three dimensions and to parallel
execution, an earlier two-dimensional finite element solver developed within
the same line of work by Gómez~\cite{gomez2019fem}.

\chapter{Scope}

The scope of the project comprises the following tasks:

\begin{itemize}
    \item The mathematical formulation of the steady and time-dependent
    advection--diffusion--reaction problems in three dimensions, including
    their weak (variational) form and the treatment of Dirichlet, Neumann and
    Robin boundary conditions.

    \item The finite element discretisation of these problems on unstructured
    tetrahedral meshes, using Lagrange elements, and the assembly of the
    resulting discrete residual and Jacobian.

    \item The import and parallel distribution of three-dimensional meshes
    generated with Gmsh, together with the identification of the different
    parts of the boundary to which different boundary conditions are applied.

    \item The implementation of the simulator in the C language on top of
    PETSc, using its \texttt{DMPlex} mesh abstraction, its \texttt{PetscFE}
    finite element support, and its \texttt{KSP}, \texttt{SNES} and
    \texttt{TS} solver frameworks.

    \item The design of a specification layer, based on the Lua language, that
    lets the user describe a problem ---coefficients, source terms and the
    type and data of each boundary condition--- without modifying or
    recompiling the C code.

    \item The validation of the simulator against problems with a known
    analytical solution and against reference solutions computed with the
    FreeFEM package~\cite{freefem-manual} on the same geometries, together with
    the assessment of its parallel behaviour by comparing sequential and
    parallel executions.
\end{itemize}

\section*{Out of Scope}

The following aspects are explicitly considered outside the scope of this project:

\begin{itemize}
    \item The generation of three-dimensional tetrahedral computational grids themselves is delegated externally to Gmsh.
    \item The study and implementation of nonlinear partial differential equations beyond the linear model equations.
    \item Execution on graphics processing units (GPUs) or other hardware acceleration platforms.
    \item Automatic or dynamic refinement techniques during the simulation process.
    \item Physical laboratory experiments or experimental testing.
\end{itemize}

\chapter{Requirements}

The simulator is subject to the following basic requirements:
\begin{itemize}
    \item It must solve the steady and time-dependent
    advection--diffusion--reaction equations on three-dimensional
    unstructured tetrahedral meshes.

    \item It must support Dirichlet, Neumann and Robin boundary conditions,
    and must allow different conditions to be imposed on different parts of the
    boundary of the same domain.

    \item The type of boundary condition applied to each part of the boundary,
    as well as all the coefficients and source terms of the problem, must be
    configurable at run time, without recompiling the code.

    \item It must run correctly in parallel on distributed-memory machines
    using the Message Passing Interface (MPI), and produce results that are
    independent of the number of processes used.

    \item It must rely exclusively on free and open-source software (PETSc,
    an MPI implementation, Gmsh and Lua), so that it is portable and
    reproducible.

    \item It must write its results in a format that can be read by standard
    scientific visualisation tools, so that the computed solution can be
    inspected and post-processed with, for example, ParaView
    \cite{paraview-guide}.
\end{itemize}
\clearpage
\chapter{Justification}

Three of the reasons that motivate this project deserve to be highlighted.

First, most problems of engineering interest are genuinely
three-dimensional. Two-dimensional models are often convenient
simplifications, but they cannot capture geometric features such as the
cooling channel of a heat sink or the true shape of a mechanical part.
Extending a finite element solver to three dimensions, and doing so in a way
that scales to realistic problem sizes, is therefore a natural and valuable
step.

Second, the size of three-dimensional problems makes parallel computing
essential. A modern open-source stack ---with PETSc \cite{petsc-user-ref} at
its core--- makes it possible to build scalable parallel solvers without
reimplementing the underlying linear algebra, mesh management or Krylov
methods \cite{saad2003iterative} from scratch, and this project takes
advantage of exactly that infrastructure.

Third, separating the description of a problem from the code that solves it

has a clear practical value. By expressing the coefficients, source terms and

boundary conditions in a small Lua specification, the same compiled simulator

can be applied to many different problems, which makes it a far more useful

tool, both in an engineering and in an educational setting.

\chapter{Schedule}

The project is structured into sequential and parallel tasks designed to fulfill the requirements and scope of the 3D finite element numerical simulator. Table~\ref{tab:schedule} details the task identification, dependencies, and descriptions.

\begin{table}[htbp]
\centering
\caption{Project Schedule}
\label{tab:schedule}
\resizebox{\textwidth}{!}{%
\begin{tabular}{|c|l|c|p{8cm}|}
\hline
\textbf{Task Code} & \textbf{Identification} & \textbf{Dependencies} & \textbf{Description} \\ \hline
\textbf{A} & Project Charter & — & Definition of the project scope, requirements, object, and schedule. \\ \hline
\textbf{B} & Mathematical Foundations Study & — & Study of the 3D boundary value problem, weak formulation, and finite element discretisation on tetrahedra. \\ \hline
\textbf{C} & PETSc \& Infrastructure Setup & B & Initial implementation of distributed data structures (\texttt{Vec}, \texttt{Mat}, \texttt{KSP}) and setup of PETSc environment. \\ \hline
\textbf{D} & 3D Domain \& Gmsh Integration & C & Implementation of the mesh import pipeline using Gmsh (\texttt{DMPlex}), parallel mesh distribution (\texttt{METIS}/\texttt{PT-Scotch}), and boundary tagging. \\ \hline
\textbf{E} & Lua Specification Layer & C & Implementation of the embedded Lua scripting interpreter to configure problem data and boundary conditions dynamically. \\ \hline
\textbf{F} & FEM Operator Assembly & D, E & Discretisation and assembly of steady and time-dependent residual and Jacobian terms using \texttt{PetscFE} and \texttt{PetscDS}. \\ \hline
\textbf{G} & Linear \& Transient Solver Integration & F & Integration of Krylov subspace methods, multigrid preconditioning (\texttt{gamg}), and \texttt{TS} time-marching schemes (Euler, Crank-Nicolson). \\ \hline
\textbf{H} & Code Verification via MMS \& Shell & G & Code verification using the Method of Manufactured Solutions and the exact analytical solution of the spherical shell. \\ \hline
\textbf{I} & Numerical Validation \& Cross-Validation & H & Numerical validation against reference solutions computed with FreeFEM on re-entrant corner domains (Fichera). \\ \hline
\textbf{J} & Mesh Refinement \& Parallel Scaling Study & I & Performance evaluation including uniform/adaptive mesh refinement (\texttt{Mmg}/\texttt{ParMmg}) and strong parallel scaling analysis. \\ \hline
\textbf{K} & Post-processing \& Visualization & G & Exporting results to \texttt{.vtu} and HDF5 formats and generating 3D renders and animations using ParaView. \\ \hline
\textbf{L} & Economic Budget \& Impact Analysis & J & Estimation of computational and personal costs, including environmental/social considerations. \\ \hline
\textbf{M} & Final Report Writing & — & Documentation of the full thesis report, mathematical formulation, code listings, and results. \\ \hline
\textbf{N} & Oral Defense Preparation & M & Preparation of presentation slides and materials for the final thesis defense. \\ \hline
\end{tabular}%
}

\end{table}
\begin{figure}[!ht]

\begin{center}
\includegraphics[width=1\textwidth]{figures/Grantt.png}
\end{center}
\caption{Project Gantt Chart (Source: Created by author)}
\label{fig:grantt}
\end{figure}

\end{document}
